% =====================================================================
% Design and Implementation of an Intelligent sEMG-Based Human-Machine
% Interface for Multi-DoF Robotic Arm Control
%
% Author: Muhtasim Al Ahsan
% Department of CSE, State University of Bangladesh
%
% Compile with:  pdflatex main && bibtex main && pdflatex main && pdflatex main
% Or upload to Overleaf and just press the green button.
% =====================================================================

\documentclass[conference]{IEEEtran}
\IEEEoverridecommandlockouts

\usepackage{cite}
\usepackage{amsmath,amssymb,amsfonts}
\usepackage{algorithmic}
\usepackage{graphicx}
\usepackage{textcomp}
\usepackage{xcolor}
\usepackage{booktabs}
\usepackage{url}
\usepackage{hyperref}

% TODO marker — search for "TODO" in the source to find what's left.
\newcommand{\todo}[1]{\textcolor{red}{\textbf{[TODO: #1]}}}

\def\BibTeX{{\rm B\kern-.05em{\sc i\kern-.025em b}\kern-.08em
    T\kern-.1667em\lower.7ex\hbox{E}\kern-.125emX}}

% =====================================================================
\begin{document}

\title{Design and Implementation of an Intelligent sEMG-Based\\
Human--Machine Interface for Multi-DoF Robotic Arm Control}

\author{
\IEEEauthorblockN{Muhtasim Al Ahsan}
\IEEEauthorblockA{Department of Computer Science \& Engineering\\
State University of Bangladesh, Dhaka, Bangladesh\\
muhtasim.ahsan74@gmail.com}
}

\maketitle

% --------------------------------------------------------------------- %
\begin{abstract}
Surface electromyography (sEMG) offers a non-invasive route to muscle-driven
human--machine interfaces, but commercial systems remain prohibitively
expensive for low-resource settings. This paper presents a complete, open
sEMG pipeline for multi-degree-of-freedom (DoF) robotic arm control,
combining a sub-\$25 analog front-end with a software stack covering
preprocessing, feature extraction, machine-learned intent decoding, and a
mode-based control state machine. The system is evaluated on
\todo{N participants, M sessions} and achieves \todo{accuracy} classification
accuracy with end-to-end latency below \todo{X}\,ms. We release the hardware
schematics, software, and a synthetic-data emulator to support
reproducibility and future research.
\end{abstract}

\begin{IEEEkeywords}
surface electromyography, sEMG, prosthetic control, multi-DoF, machine
learning, low-cost biomedical instrumentation, human--machine interface
\end{IEEEkeywords}

% --------------------------------------------------------------------- %
\section{Introduction}
\label{sec:intro}

\todo{Open with a concrete motivation: upper-limb amputation prevalence in
South Asia, cost of commercial myoelectric prostheses, the affordability
gap.}

Surface electromyography measures the electrical activity of skeletal
muscle through the skin. Because the recorded signal reflects voluntary
motor commands, sEMG has long been used to drive prosthetic limbs and
rehabilitation robots \cite{oskoei2007myoelectric, atzori2014review}.
Modern myoelectric prostheses, however, rely on commercial acquisition
hardware whose cost places them out of reach for the majority of patients
in low-resource settings.

This work asks how far an undergraduate-grade, sub-\$25 sEMG front-end can
be pushed when paired with a careful software pipeline. Specifically, we
target real-time, multi-DoF control of a low-cost servo-driven arm from
only two sEMG channels.

\textbf{Contributions.} This paper makes the following contributions:
\begin{enumerate}
    \item An open, low-cost (\textless\,\$25) analog sEMG front-end based on the AD620 instrumentation amplifier, characterised against \todo{reference sensor / signal-quality baseline}.
    \item A reproducible Python pipeline covering preprocessing, time- and frequency-domain feature extraction, and Random Forest / 1D-CNN classifiers.
    \item A mode-based control state machine that maps two EMG channels to six robotic DoFs without flat $K$-class classification.
    \item A PyQt6 graphical interface for cue-based training, real-time signal monitoring, and virtual-arm preview, supporting both hardware acquisition and a synthetic-signal emulator.
    \item Cross-session evaluation on \todo{N} healthy volunteers and a public release of code, schematics, and dataset.
\end{enumerate}

% --------------------------------------------------------------------- %
\section{Related Work}
\label{sec:related}

\subsection{sEMG-based prosthetic control}
\todo{Summarise pattern-recognition control: Hudgins TD features
\cite{hudgins1993new}, Englehart real-time scheme \cite{englehart2003robust},
NinaPro dataset, MyoArmband studies.}

\subsection{Deep learning for EMG}
\todo{Atzori et al.\ CNN baselines on NinaPro \cite{atzori2016deep};
Chen et al.\ gesture recognition \cite{chen2020deep}; transformer-based
work where relevant.}

\subsection{Mode-based and sequential control}
\todo{Co-contraction switching schemes used in commercial prostheses;
limitations of flat $K$-class classifiers when DoFs $> K$.}

\subsection{Low-cost EMG systems}
\todo{MyoWare \cite{myoware}, OpenBCI \cite{openbci}, BITalino
\cite{bitalino}; gap: most reported low-cost systems are not characterised
against a clinical reference.}

% --------------------------------------------------------------------- %
\section{Hardware}
\label{sec:hardware}

\subsection{Analog front-end}
Each channel uses a single AD620 instrumentation amplifier with a gain of
\todo{G} set by an external resistor, followed by an active bandpass stage
(20--450~Hz) and a 50~Hz twin-T notch to suppress mains coupling. A 9~V
battery and a TL072-based virtual-ground rail isolate the analog stage
from mains noise.

\todo{Insert circuit schematic (Fig.~\ref{fig:schematic}).}

\subsection{Acquisition}
A single Arduino Nano samples two channels at \todo{1\,kHz raw / 200\,Hz
envelope} through its 10-bit ADC and forwards integer samples over USB
serial at 115\,200~baud, formatted as CSV lines.

\subsection{Bill of materials}
\todo{Insert BOM table (Table~\ref{tab:bom}).}

% --------------------------------------------------------------------- %
\section{Software Pipeline}
\label{sec:software}

\subsection{Preprocessing}
For each window of incoming samples we apply, in order: per-channel DC
removal, a 4th-order Butterworth bandpass between 20~Hz and 450~Hz, and an
IIR notch at 50~Hz with quality factor $Q=30$. Streaming variants
(\texttt{LiveBandpass}, \texttt{LiveNotch}) retain filter state between
chunks to avoid edge transients.

\subsection{Envelope and segmentation}
The conditioned signal is rectified and smoothed by a 100~ms moving-average
window. We then form fixed-length overlapping windows of $T = $
\todo{$T_{\mathrm{win}}$}~ms with a hop of \todo{$T_{\mathrm{hop}}$}~ms.

\subsection{Feature extraction}
We compute the Hudgins time-domain (TD) feature set per channel:
mean absolute value (MAV), root mean square (RMS), waveform length (WL),
zero crossings (ZC) and slope-sign changes (SSC). We additionally compute
two frequency-domain descriptors, mean frequency (MNF) and median
frequency (MDF), giving a per-window feature vector of length $7C$ for
$C$ channels.

\subsection{Classification models}
Three classifiers are evaluated:
\begin{itemize}
    \item \textbf{Random Forest} with 300 trees and balanced class weights;
    \item \textbf{Linear SVM} with squared-hinge loss;
    \item \textbf{1D Convolutional Neural Network} operating directly on raw windowed signals (Section~\ref{sec:cnn}).
\end{itemize}

\subsubsection{1D CNN architecture}
\label{sec:cnn}
\todo{Describe Conv1D layer sizes, BatchNorm, max-pool, global average
pool, dropout, and softmax head. Cite the corresponding code block in the
release.}

\subsection{Mode-based control state machine}
With only two channels we cannot directly produce six DoF commands. We
instead route four classifier outputs ($\{\text{REST}, \text{CH1},
\text{CH2}, \text{BOTH}\}$) through a finite-state controller with three
modes ($\text{GRIP}$, $\text{WRIST}$, $\text{ARM}$). Sustained
co-contraction (\text{BOTH}) above a threshold for $N$ consecutive
windows advances the mode, with a cooldown period to suppress
double-switches. Within a mode, the active channel and contraction
strength drive the corresponding pair of DoFs.

% --------------------------------------------------------------------- %
\section{Experimental Setup}
\label{sec:experiments}

\subsection{Participants and protocol}
\todo{Number of volunteers; age; informed consent; ethics statement.}
Each participant completed \todo{$K$} sessions on separate days. Within a
session, the cue-based GUI presented \todo{$R$} trials per class
($\{$REST, CH1, CH2, BOTH$\}$), with a 3~s action phase and a 2~s rest
phase. All windows recorded during action phases were labelled
automatically with the cued class.

\subsection{Metrics}
\begin{itemize}
    \item \textbf{Classification accuracy} on held-out windows (within-session, cross-session, cross-user).
    \item \textbf{End-to-end latency} measured from a synthetic spike injected on the Arduino to the time at which the controller emits a command.
    \item \textbf{False-activation rate} per minute during sustained rest.
\end{itemize}

% --------------------------------------------------------------------- %
\section{Results}
\label{sec:results}

\subsection{Within-session classification}
\todo{Report Random Forest / SVM / CNN accuracy. Include confusion matrix (Fig.~\ref{fig:confusion}).}

\subsection{Cross-session and cross-user generalisation}
\todo{Train on session 1, test on sessions 2--$K$. Train on $N-1$ users, test on the held-out user. Discuss the drop.}

\subsection{Latency and stability}
\todo{Report mean$\pm$std end-to-end latency. Report false activations per minute at rest.}

\subsection{Ablations}
\todo{What happens if you drop ZC/SSC? Drop the notch filter? Reduce window size?}

% --------------------------------------------------------------------- %
\section{Discussion}
\label{sec:discussion}

\todo{Compare to MyoWare/OpenBCI benchmarks. Discuss the cost--accuracy
trade-off. Address signal drift, electrode shift, and fatigue. Note that
sEMG decodes muscle activation, not abstract thought, and that this is
identical to how commercial myoelectric prostheses operate.}

% --------------------------------------------------------------------- %
\section{Limitations and Future Work}
\label{sec:limitations}

This work has several limitations.
First, the analog front-end is single-supply and battery-only, which
limits common-mode rejection compared to clinical-grade differential
preamplifiers.
Second, two channels of envelope sEMG cannot support fully
proportional 6-DoF control; the mode-based scheme trades expressivity for
robustness.
Third, our evaluation is limited to able-bodied participants. Validation
with upper-limb amputees, which would require institutional ethics
review, remains future work.
A natural extension is to add an EEG channel for higher-level mode
switching, which would complement rather than replace the sEMG channels.

% --------------------------------------------------------------------- %
\section{Conclusion}
\label{sec:conclusion}

We presented a complete, low-cost sEMG-based interface for multi-DoF
robotic arm control. The system reaches \todo{X}\% within-session
accuracy and \todo{Y}\,ms end-to-end latency at a hardware cost below
\todo{Z} USD, and is released as open hardware and open source to
support reproducibility and reuse in low-resource settings.

% --------------------------------------------------------------------- %
\section*{Acknowledgments}
\todo{Thank your supervisor at SUB CSE, the Rufaida BioMeds team for
prototyping support, and any volunteers.}

% --------------------------------------------------------------------- %
\bibliographystyle{IEEEtran}
\begin{thebibliography}{99}

\bibitem{oskoei2007myoelectric}
M.~A.\ Oskoei and H.\ Hu, ``Myoelectric control systems---A survey,''
\emph{Biomed.\ Signal Process.\ Control}, vol.~2, no.~4, pp.~275--294, 2007.

\bibitem{atzori2014review}
M.\ Atzori \emph{et al.}, ``Electromyography data for non-invasive
naturally-controlled robotic hand prostheses,'' \emph{Sci.\ Data}, vol.~1,
2014.

\bibitem{hudgins1993new}
B.\ Hudgins, P.\ Parker, and R.\ N.\ Scott, ``A new strategy for
multifunction myoelectric control,'' \emph{IEEE Trans.\ Biomed.\ Eng.},
vol.~40, no.~1, pp.~82--94, Jan.\ 1993.

\bibitem{englehart2003robust}
K.\ Englehart and B.\ Hudgins, ``A robust, real-time control scheme for
multifunction myoelectric control,'' \emph{IEEE Trans.\ Biomed.\ Eng.},
vol.~50, no.~7, pp.~848--854, Jul.\ 2003.

\bibitem{atzori2016deep}
M.\ Atzori, M.\ Cognolato, and H.\ M{\"u}ller, ``Deep learning with
convolutional neural networks applied to electromyography data: A
resource for the classification of movements for prosthetic hands,''
\emph{Front.\ Neurorobot.}, vol.~10, 2016.

\bibitem{chen2020deep}
X.\ Chen \emph{et al.}, ``Deep learning for EMG-based hand gesture
recognition: A review,'' \emph{IEEE Access}, vol.~8, pp.~225--241, 2020.

\bibitem{myoware}
Advancer Technologies, ``MyoWare Muscle Sensor 2.0,'' product
documentation, 2022.

\bibitem{openbci}
OpenBCI, ``Cyton biosensing board,'' product documentation, 2024.

\bibitem{bitalino}
PLUX Biosignals, ``BITalino (r)evolution,'' product documentation, 2024.

\todo{Add 10--15 more references during the literature review phase.}

\end{thebibliography}

\end{document}
